\section{Conclusion}
\label{sec:conclusion}

AI disputes arrive in court wearing metaphors. A chatbot ``promises.'' A companion ``cares.'' An agent ``decides.'' A model ``speaks.'' The metaphors are useful descriptions of an encounter, but they do not identify the legal parties, allocate authority, or make a judgment enforceable. Courts need not wait for a final philosophy of machine mind to perform those tasks. They need a minimum structure that remains stable while technology and sectoral regulation change.

The first element is a stable adjudicative object. AI law can be a coherent field without a universal technical definition or comprehensive code. Its recurring work is to organize relations among people, organizations, automated systems, and public institutions. Regulatory definitions may therefore remain purpose specific. A medical-device rule, a companion-chatbot statute, a procurement condition, and a copyright doctrine can draw different boundaries. The invariant lies beneath them: a functional classification cannot silently rewrite the cast of legal subjects.

The second element is status without a Turing test. Operational autonomy does not generate artificial responsibility personality. A system can plan, adapt, initiate tool calls, and produce language that appears intentional. Those capabilities can transform foreseeable risk and the precautions reasonable deployment requires. They do not by themselves supply representation, assets, claims, duties, procedural capacity, sanctions, or enforcement. Governing positive law may one day construct a limited AI-related legal platform, just as it constructs corporations and other juridical entities. Until then, every normative position necessary to judgment must terminate in a recognized natural or juridical person. The system remains an automatic instrumentality capable of supplying conduct, causation, and evidence without becoming an assetless defendant.

The third element is a reconstructed control chain. No human may select the harmful sentence at the instant it appears, yet organizations select and maintain the mechanism that generates it. The deployment-control cascade makes that control visible. It traces objectives and resources downward through provider, model, harness, platform, tools, and user environment; it traces evaluations, complaints, incidents, and other feedback upward; and it identifies who can override, roll back, restrict, or stop operation. Its variables translate into familiar law: authority can bear on duty and attribution, feedback on notice, process models on foreseeability, stopping power on feasible precaution, and records on proof. The cascade does not decide liability. It ensures that probabilistic operation does not erase the facts from which liability is decided.

The fourth element is responsibility in two layers. Externally, the enterprise that makes or maintains the deployment in its activity should be the claimant-facing responsibility center. That rule does not guarantee recovery and does not turn every hallucination into a tort. The claimant must prove a recognized injury, breach or defect, causation, and the other elements and defenses supplied by governing law. The rule does prevent the enterprise from requiring an injured person to name the responsible engineer before discovery or from treating ``AI-generated'' as the end of attribution. For high-impact deployments, the enterprise should also produce the deployment and incident record needed to locate the real control path.

Internally, a mandatory responsibility charter should fix authority before the loss. It should name risk owners, record deployment approval and dissent, protect escalation, confer usable stop and rollback power, preserve incident and configuration evidence, and specify indemnification and recourse. A worker who documents a danger and uses the authorized channel should not occupy the same position as a manager who suppresses the report or a developer who secretly removes a constraint. The enterprise faces the public; the charter protects the reporter and pursues the concealer. That alignment gives safety information somewhere to go and legal responsibility somewhere effective to land.

The framework is deliberately reciprocal. It permits different capabilities for public, professional, research, and government deployments while requiring consequential boundaries to be proportionate and reviewable. It permits AI to extend care and administration while refusing to let automation alone extinguish an existing duty. It permits government to choose mission requirements and suppliers while distinguishing procurement from retaliation or unsupported coercion. It treats provider constitutions and state directives alike: each can organize internal conduct, and each remains subject to external law.

The Hangzhou chatbot, Character.AI, a clinical model, and Claude in a defense environment present profoundly different facts. They can nevertheless be adjudicated through the same four moves: identify the existing legal subjects; classify the system for the doctrine actually invoked; reconstruct the deployment's control and feedback; and allocate external and internal responsibility without inventing a machine person. That is enough jurisprudence to begin deciding well. It is also the structure on which more specialized AI law can safely build.
